\newglossaryentry{QGIS}
{
    name=QGIS,
    description={An open-source geographic information system.}
}
%\newacronym{IC}{IC}{Item Count}
\newacronym{SE}{SE}{Siemens Energy}
\newacronym{CM}{CM}{Condition Monitoring}
\newacronym{TSA}{TSA}{Time Synchronous Averaging}
\newacronym{ANC}{ANC}{Adaptive Noise Cancellation}
\newacronym{SANC}{SANC}{Self-Adaptive Noise Cancellation}
\newacronym{DRS}{DRS}{Discrete and Random Signal Separation}
\newacronym{SES}{SES}{Squared Envelope Spectrum}
\newacronym{RMSE}{RMSE}{Root Mean Square Error}
\newacronym{RMS}{RMS}{Root Mean Square}
\newacronym{TP}{TP}{True Positive}
\newacronym{TN}{TN}{True Negative}
\newacronym{FP}{FP}{False Positive}
\newacronym{FN}{FN}{False Negative}

\newglossaryentry{filterbank}{
    name={Filterbank},
    description={Anordnung mehrerer Filter, die ein Signal in verschiedene
    Frequenzbänder zerlegen
    \cite[S.~110, Abschn.~2.1--2.2]{antoniFastComputationKurtogram2007}.
    \hyperlink{begriff:filterbank}{Zur Textstelle in Abschnitt~\ref*{sec:bandauswahl}}.}
}
\newglossaryentry{dyadisches-raster}{
    name={Dyadisches Raster},
    text={dyadisches Raster},
    description={Raster aus Frequenzbändern, deren Anzahl sich bei jedem
    Zerlegungsschritt verdoppelt, während sich ihre Bandbreite halbiert
    \cite[S.~110--111, Abschn.~2.2]{antoniFastComputationKurtogram2007}.
    \hyperlink{begriff:dyadisches-raster}{Zur Textstelle in Abschnitt~\ref*{sec:bandauswahl}}.}
}
\newglossaryentry{drittel-binaerbaum}{
    name={1/3-Binärbaum},
    sort={Drittel-Binaerbaum},
    description={Filterbankstruktur des \emph{Fast Kurtogram}, die die binäre
    Zerlegung um eine Dreiteilung der Teilbänder ergänzt. Dadurch entstehen
    zusätzliche Zwischenebenen im Raster aus Frequenz und Bandbreite
    \cite[S.~114, Abschn.~2.5, Abb.~5]{antoniFastComputationKurtogram2007}.
    \hyperlink{begriff:drittel-binaerbaum}{Zur Textstelle in Abschnitt~\ref*{sec:bandauswahl}}.}
}
\newglossaryentry{stft}{
    name={Kurzzeit-Fourier-Transformation (STFT)},
    text={Kurzzeit-Fourier-Transformation},
    sort={Kurzzeit-Fourier-Transformation},
    description={Die \emph{Short-Time Fourier Transform} berechnet die
    Fourier-Transformation eines Signals in zeitlich verschobenen
    Analysefenstern. Sie beschreibt damit den Frequenzinhalt in Abhängigkeit
    von der Zeitposition des Fensters
    \cite[S.~294, Abschn.~5.1, Gl.~(36)]{antoniSpectralKurtosisUseful2006}.
    \hyperlink{begriff:stft}{Zur Textstelle in Abschnitt~\ref*{sec:bandauswahl}}.}
}
